\section{Seção com Matemática}

Texto com matemática inline: $e^{i\pi}+1=0$, $\alpha+\beta=\gamma$ e $x^2+y^2=z^2$.

\subsection{Ambiente equation}

\begin{equation}
  \label{eq:integral}
  \int_{0}^{\infty} e^{-x^2}\,dx = \frac{\sqrt{\pi}}{2}.
\end{equation}

\subsection{Ambiente align}

\begin{align}
  (a+b)^2 &= a^2 + 2ab + b^2, \\
  (a-b)^2 &= a^2 - 2ab + b^2.
\end{align}

\subsection{Ambientes Matemáticos}

Os ambientes seguem o estilo visual do template e são compatíveis com \verb|\cref|.

\begin{definicao}[Transformada Z]
  \label{def:transformada-z}
  A Transformada Z bilateral de um sinal discreto $x[n]$ é definida como
  \[
    X(z) = \sum_{n=-\infty}^{\infty} x[n]\, z^{-n}, \quad z \in \mathbb{C}.
  \]
\end{definicao}

\begin{teorema}[Deslocamento no Tempo]
  \label{teo:deslocamento}
  Se $X(z) = \mathcal{Z}\{x[n]\}$, então para qualquer inteiro $k$:
  \[
    \mathcal{Z}\{x[n-k]\} = z^{-k} X(z).
  \]
\end{teorema}

\begin{proof}
  Substituindo diretamente na \cref{def:transformada-z}:
  \[
    \sum_{n=-\infty}^{\infty} x[n-k]\,z^{-n}
    = z^{-k} \sum_{m=-\infty}^{\infty} x[m]\,z^{-m}
    = z^{-k} X(z). \qedhere
  \]
\end{proof}

\begin{corolario}
  O atraso de um sample corresponde a multiplicar por $z^{-1}$ no domínio Z.
\end{corolario}

\begin{exemplo}
  Para $x[n] = \delta[n]$ (impulso unitário), tem-se $X(z) = 1$.
  Pelo \cref{teo:deslocamento}, $\mathcal{Z}\{\delta[n-1]\} = z^{-1}$.
\end{exemplo}

\begin{observacao}
  A região de convergência (ROC) deve ser sempre especificada,
  pois sequências distintas podem ter a mesma expressão $X(z)$.
\end{observacao}
